\providecommand{\frA}{\mathfrak A}
\providecommand{\frB}{\mathfrak B}
\providecommand{\frC}{\mathfrak C}
\providecommand{\frR}{\mathfrak R}
\providecommand{\frS}{\mathfrak S}
\providecommand{\frT}{\mathfrak T}
\providecommand{\frM}{\mathfrak M}
\providecommand{\frG}{\mathfrak G}
\providecommand{\frako}{\mathfrak o}

\section*{40. Algèbres non commutatives}

\begin{center}
\emph{Math. Zs. 37 (1933), p. 514--541}
\end{center}

\begin{center}
\textbf{Algèbre non commutative.}\\[0.5em]
Par\\[0.25em]
Emmy Noether à Goettingen.
\end{center}

Les théorèmes principaux de l'algèbre commutative sont, comme on sait, contenus dans la théorie de Galois; celle-ci est précédée par la théorie des corps de rupture et des corps de décomposition, c'est-à-dire des corps qui suffisent pour qu'un polynôme donné fasse apparaître un facteur linéaire, respectivement se décompose entièrement en facteurs linéaires.

Je développe ici les parties correspondantes de l'algèbre dans le non-commutatif, en particulier dans l'hypercomplexe. Et je travaille alors par principe avec des méthodes non commutatives, avec la représentation dans des corps gauches. À la fin je montre comment les théorèmes mentionnés ci-dessus de l'algèbre commutative peuvent se fonder d'une manière tout à fait parallèle par représentation dans des corps commutatifs.

La théorie des représentations sur laquelle repose ce travail --- précédée d'une brève théorie des endomorphismes (\S 1) --- prolonge celle qui est fondée sur les modules de représentation.\footnote{E. Noether, Hyperkomplexe Grossen und Darstellungstheorie, Math. Zeitschr. 30 (1929), p. 641--692, cité ci-dessous comme Darstellungstheorie. Voir la reproduction chez van der Waerden, Moderne Algebra II.} En particulier, je considère, à côté de la représentation directe, la représentation réciproque --- chacune pouvant se ramener à l'autre ---, désormais produite au moyen du module de représentation réciproque. L'avantage est que ce module --- qui est donc un bimodule --- peut aussi être regardé comme un module unilatéral sur un anneau d'extension (\S 2). La représentation dans des corps gauches se ramène ainsi à la théorie des idéaux de l'anneau d'extension; les classes de représentations irréductibles correspondent aux classes d'idéaux irréductibles de cet anneau, en exacte analogie avec les faits qui valent pour le système lui-même lorsqu'on représente un système hypercomplexe sur son corps commutatif de coefficients (\S 3 et 4).

On en déduit les théorèmes structuraux sur les anneaux de matrices sur des corps gauches, grâce à la remarque (\S 5) selon laquelle \emph{tout sous-anneau fournit une représentation par ce corps gauche, ou une représentation réciproque par le corps gauche anti-isomorphe}. Ce corps gauche anti-isomorphe constitue un premier analogue du corps de scission minimal et du corps de décomposition dans le cas commutatif: il réalise toutes les représentations réciproques de degré 1 et, lorsque le rang sur le centre --- qui n'avait pas été supposé fini jusque-là --- est fini, une décomposition complète en facteurs directs de rang 1. On obtient ainsi la théorie de Galois pour les corps gauches (\S 6), ainsi que le fait (\S 7) que des algèbres à division anti-isomorphes (corps gauches de rang fini sur leur centre) déterminent des classes inverses dans le groupe de R. Brauer des classes d'algèbres.

Une seconde démonstration de la théorie de Galois, valable plus généralement pour les systèmes simples (et présentée plus tôt dans le texte), découle presque immédiatement d'un théorème sur les sous-anneaux commutant entre eux, lui-même rattaché presque directement à l'étude préliminaire des endomorphismes.

Jusqu'ici, les considérations sont entièrement non commutatives. La question des corps de décomposition \emph{commutatifs} est néanmoins elle aussi traitée par des méthodes non commutatives, au moyen de la représentation du corps de décomposition par l'algèbre à division, conformément à la remarque préparée au \S 5 (\S 7). L'article se conclut par le passage au cas commutatif mentionné plus haut (\S 8) et par la théorie des corps de décomposition des systèmes arbitraires (\S 9); on obtient en même temps le lien avec la théorie usuelle des représentations sur un corps commutatif.

R. Brauer a fondé la théorie des corps de décomposition sur cette théorie de la représentation dans le commutatif --- et sur les systèmes de facteurs ``irrationnels''.\footnote{Voir notre note commune: Uber minimale Zerfallungskorper irreduzibler Darstellungen. Sitz. Ber. d. Preuss. Ak. d. Wiss. 1927, p. 221--228. (La remarque de la p. 222 contient un ``rapport justificatif''. Les théorèmes principaux sont nés indépendamment et presque simultanément.) Voir aussi R. Brauer, Uber Systeme hyperkomplexer Zahlen, \S 3. Math. Zeitschr. 30 (1929), p. 79--107. A. A. Albert a retrouvé plus tard ces théorèmes de façon indépendante.} À cause de cette fondation commutative il doit --- tout comme Albert plus tard --- supposer que le centre est un corps parfait, restriction inutile d'après la fondation non commutative. Avec la même restriction, et également avec la théorie de la représentation dans le commutatif, R. Brauer et K. Shoda ont --- après avoir pris connaissance de ma théorie de Galois pour les corps gauches --- développé davantage la théorie: R. Brauer donna le théorème mentionné ci-dessus sur les sous-anneaux permutables, K. Shoda indépendamment la théorie complète aussi pour les systèmes semi-simples.\footnote{R. Brauer, Uber die algebraische Struktur von Schiefkorpern, \S 2. Journ. f. Math. 166 (1932), p. 241--252. K. Shoda, Uber die galoissche Theorie der halbeinfachen hyperkomplexen Systeme. Math. Ann. 107 (1932), p. 252--258. Ces résultats avaient aussi été développés par J. Levitzki.}

Enfin il faut mentionner que, pour le cas des corps gauches, on trouve une brève présentation chez van der Waerden II (\S 128), à la suite de mon cours de l'été 1928, où j'avais parlé pour la première fois de ces questions. La présentation de van der Waerden apporte une série de simplifications par rapport à ce cours; je les ai reprises et prolongées en partie dans un second cours (hiver 1929/30),\footnote{Le premier cours a été élaboré par G. Kothe, le second par M. Deuring.} en partie seulement ici. En particulier, le transfert de la méthode de preuve hypercomplexe de la théorie de Galois au commutatif (\S 8) ne date que du second cours, où la preuve non commutative (\S 6, 3.) avait été essentiellement simplifiée par rapport au premier cours. Le théorème sur les sous-anneaux permutables (\S 5) et la seconde méthode de preuve de la théorie de Galois non commutative qui s'appuie sur lui (\S 6, 1. et 2.) n'ont été ajoutés qu'après le second cours, après connaissance de R. Brauer (note 3) et de van der Waerden \S 128; dans \S 5, 3., les hypothèses de finitude sont toutefois beaucoup plus faibles que là.

Certains raisonnements du premier cours --- méthode de formation d'intersections d'idéaux de différences --- ont, bien que plus compliqués, l'avantage de se transposer aux systèmes de rang infini et aux systèmes entiers.\footnote{La méthode est reproduite pour l'essentiel chez G. Kothe, Schiefkorper unendlichen Ranges uber dem Zentrum, \S 5, Math. Ann. 105 (1931), p. 15--39. Voir aussi la note 10. Cette méthode d'intersection est maintenant remplacée par le fait presque trivial (\S 4, 1.) --- remarqué pour la première fois par van der Waerden --- que les idéaux bilatères appartiennent aux modules invariants. Tout peut ainsi se fonder par la décomposition plus simple en sommes directes, ce qui échoue cependant dans le cas du rang infini et dans le cas entier.} Pour les systèmes commutatifs entiers, on parvient ainsi au lien entre différentiation idéale et différente,\footnote{Voir la conférence de Prague, Jahresber. d. Deutsch. Math. Ver. 39 (1930), p. 17 (pagination oblique).} sur lequel je reviendrai à l'occasion plus exactement.

La théorie des corps de décomposition trouve son application principale dans la théorie des produits croisés et de leurs systèmes de facteurs, qui constituent à leur tour le fondement d'applications arithmétiques; il ne sera plus question ici de cela.\footnote{La théorie des produits croisés est développée dans le second cours; elle est reproduite --- avec de petites modifications, pour ne pas devoir supposer la théorie donnée ici --- chez H. Hasse: Theory of cyclic algebras over an algebraic numberfield, chap. II. Transactions of the Amer. Math. Soc. 34 (1932), p. 171--214. Une présentation plus étroitement rattachée au cours paraîtra dans un rapport de M. Deuring dans les ``Ergebnisse der Mathematik''.}

\subsection*{\S 1. Endomorphismes, modules et bimodules}

La théorie de la représentation fondée sur les modules de représentation repose sur la théorie de l'anneau des endomorphismes des groupes abéliens, avec ou sans opérateurs. Les raisonnements implicitement sous-jacents, donc les relations entre application et lois de calcul, seront formulés dans quelques propositions simples afin d'éviter les répétitions (voir aussi \S 2).

\paragraph{1. Application multiplicative. Loi associative.}
Soit d'abord \(\frG\) un groupe sans opérateurs, et soit \(\frA\) son domaine absolu des endomorphismes, c'est-à-dire le système de tous les homomorphismes de \(\frG\) dans lui-même. \(\frA\) est fermé multiplicativement, puisque le produit \( \sigma\tau \) est défini par
\[
        g(\sigma\tau)=(g\sigma)\tau,
        \qquad g\in\frG,\quad \sigma,\tau\in\frA                         \tag{1}
\]
et satisfait évidemment à la loi associative.

\(\frG\) devient, comme on sait, un groupe avec opérateurs lorsqu'est donnée une famille \(\frB\) de symboles \(\Theta,H,\ldots\), de telle sorte que les combinaisons \(g\Theta,gH,\ldots\) fournissent des éléments de \(\frG\) définis de façon univoque et engendrent des endomorphismes (homomorphismes de \(\frG\) dans lui-même):
\[
        (gh)\Theta=g\Theta\cdot h\Theta .
\]
Il existe donc une application univoque (en général non biunivoque) de \(\frB\) sur un sous-ensemble \(\overline{\frB}\) du domaine absolu des endomorphismes \(\frA\). Si le domaine d'opérateurs \(\frB\) est fermé multiplicativement, l'application univoque de \(\frB\) sur la partie correspondante du domaine des endomorphismes est multiplicativement homomorphe si et seulement si la relation associative
\[
        g(\Theta H)=(g\Theta)H, \qquad g\in\frG,\ \Theta,H\in\frB,                       \tag{1a}
\]
est satisfaite. De façon correspondante, l'antihomomorphisme est défini lorsqu'il s'agit d'opérateurs à gauche.

En effet, l'application de \(\frB\) sur \(\overline{\frB}\) est donnée par \(g\Theta=g\sigma\) pour tout \(g\) de \(\frG\), donc aussi par
\[
        (g\Theta)H=(g\sigma)\tau=g(\sigma\tau),
\]
de sorte que (1a) devient nécessaire et suffisante pour l'homomorphie multiplicative. Le fait que les opérateurs à gauche engendrent un antihomomorphisme vient de ce que les endomorphismes écrits à gauche \(\sigma^\ast g,\tau^\ast\sigma^\ast g\) doivent se lire de droite à gauche: \(\tau^\ast\sigma^\ast g=g\sigma\tau\), tandis que les opérateurs se lisent toujours de gauche à droite.

\paragraph{2. Application homomorphe relativement aux opérateurs.}
Si \(\frG\) est un groupe avec opérateurs, l'homomorphie relativement aux opérateurs est définie par
\[
        (g\Theta)\sigma=(g\sigma)\Theta,                                          \tag{2}
\]
respectivement, pour les opérateurs à gauche, par
\[
        (\Theta g)\sigma=\Theta(g\sigma).                                          \tag{2*}
\]
c'est-à-dire par la commutation des deux actions ou, dans le cas des opérateurs à gauche, par la loi d'associativité mixte. En appliquant, comme en 1., le raisonnement qui identifie un domaine d'opérateurs à son image dans le domaine des endomorphismes, on obtient pour deux domaines d'opérateurs distincts:

Si deux domaines d'opérateurs \(\frB\) et \(\frC\) sont donnés, alors les éléments de \(\frC\) engendrent des \(\frB\)-endomorphismes, et simultanément ceux de \(\frB\) des \(\frC\)-endomorphismes, si et seulement si les actions des deux domaines commutent sur \(\frG\) ou, respectivement, si la loi d'associativité mixte est satisfaite:
\[
        (g\Theta)\bar H=(g\bar H)\Theta,                                                    \tag{2a}
\]
respectivement
\[
        (\Theta g)\bar H=\Theta(g\bar H).                               \tag{2a*}
\]

\paragraph{3. Modules et bimodules sur des anneaux.}
Un module à droite \(\frM\) sur un anneau \(\frR\) est un groupe abélien additif avec les éléments de \(\frR\) comme opérateurs à droite; outre la relation associative (1a), les relations distributives y sont satisfaites:
\[
        (g+h)\rho=g\rho+h\rho,\qquad
        g(\rho+\sigma)=g\rho+g\sigma.                                  \tag{3a}
\]
Il en va de même pour les modules à gauche.

Parmi les bimodules il faut distinguer deux espèces. Les modules à droite sur deux anneaux \(\frR\) et \(\frS\) s'appellent bimodules lorsque, outre les combinaisons valables séparément pour \(\frR\) et \(\frS\), on a encore
\[
        (m\rho)\sigma=(m\sigma)\rho
\]
Un groupe qui est un \(\frR\)-module à gauche et un \(\frS\)-module à droite est appelé bimodule lorsque s'ajoute la loi d'associativité mixte
\[
        \rho(m\sigma)=(\rho m)\sigma .
\]

La signification de ces combinaisons vient du fait que, dans le cas des groupes abéliens, le domaine absolu des endomorphismes forme un anneau, l'anneau absolu d'endomorphismes\footnote{Pour les groupes non abéliens il s'agit d'un anneau ``généralisé''; voir à ce sujet un travail de H. Fitting à paraître dans les Math. Annalen. [Math. Annalen 107, p. 514--542.]}. Dans le détail, le raisonnement par application donne:

Si le domaine d'opérateurs \(\frR\) d'un groupe abélien (écrit additivement) \(\frM\) est un anneau, l'application canonique de \(\frR\) dans l'anneau absolu d'endomorphismes, d'image \(\overline{\frR}\), est un homomorphisme d'anneaux si et seulement si \(\frM\) est un module à droite sur \(\frR\) -- donc si (1a) et (3a) sont satisfaites --, et un antihomomorphisme d'anneaux pour les modules à gauche.

Si \(\frM\) est un module à droite sur les anneaux \(\frR\) et \(\frS\), alors \(\frM\) est un bimodule si et seulement si l'application de \(\frR\) dans l'anneau des \(\frS\)-endomorphismes définie par son action est un homomorphisme d'anneaux et, simultanément, si l'application correspondante de \(\frS\) dans l'anneau des \(\frR\)-endomorphismes en est un. Si \(\frM\) est un module à gauche sur \(\frR\) et un module à droite sur \(\frS\), la condition de bimodule (2a*) signifie que l'application de \(\frS\) dans l'anneau des \(\frR\)-endomorphismes est un homomorphisme d'anneaux, tandis que celle de \(\frR\) dans l'anneau des \(\frS\)-endomorphismes est un antihomomorphisme d'anneaux.

Il en résulte en particulier:

1) Si \(\frR\) est un anneau unitaire, alors \(\frR\) est isomorphe à son anneau des endomorphismes comme \(\frR\)-module à gauche et anti-isomorphe à son anneau des endomorphismes comme \(\frR\)-module à droite. En effet, la loi d'associativité mixte (2a*) est satisfaite, donc l'assertion d'homomorphie s'applique. L'existence de l'élément unité donne tout endomorphisme sous la forme \(e\mapsto a\), avec \(a\in\frR\); on a donc un isomorphisme, et \(\frR\) épuise l'anneau des \(\frR\)-endomorphismes.

2) Si \(\frR\) est anneau complet de matrices sur un corps gauche \(A\), en général non commutatif,
\[
        \frR=\sum c_{ik}A=\sum A c_{ik},
\]
alors \(A\) est anti-isomorphe au corps gauche des endomorphismes des idéaux à droite simples, et isomorphe au corps gauche des endomorphismes des idéaux à gauche simples. En effet, tout idéal à droite simple est isomorphe, comme module à opérateurs, à un \(A\)-module à gauche et \(\frR\)-module à droite, et \(A\) épuise les \(\frR\)-endomorphismes; de même pour les idéaux à gauche.

\paragraph{4. Passage d'un bimodule à droite à un module unilatéral sur un anneau-produit.}
L'intérêt des bimodules à droite repose essentiellement sur ce passage.

\paragraph{Théorème.}
Si \(\frM\) est un module à droite sur un anneau \(\frT\) contenant deux sous-anneaux \(\frR\) et \(\frS\) qui commutent élément par élément, alors \(\frM\) peut aussi être regardé comme un bimodule sur \(\frR\) et \(\frS\). Réciproquement, si l'on se donne un bimodule sur \(\frR\) et \(\frS\) et qu'il existe un anneau-produit \(\frT\) dans lequel \(\frR\) et \(\frS\) commutent élément par élément, alors \(\frM\) peut aussi être regardé comme un module à droite sur \(\frT\), l'action de \(\frT\) prolongeant les actions données de \(\frR\) et de \(\frS\).

La première assertion vient de ce que, par définition, l'action de \(\frT\) définit un homomorphisme de \(\frT\) sur un sous-anneau \(\overline{\frT}\) de l'anneau absolu d'endomorphismes, tandis que les images \(\overline{\frR}\) et \(\overline{\frS}\) commutent élément par élément. C'est précisément la relation (2a) qui caractérise le bimodule.

Réciproquement, si \(\frM\) est donné comme \(\frR,\frS\)-bimodule (module à droite), les actions de \(\frR\) et de \(\frS\) définissent des homomorphismes vers des sous-anneaux \(\overline{\frR},\overline{\frS}\) de l'anneau absolu d'endomorphismes, dont les éléments commutent deux à deux. Dans l'anneau absolu d'endomorphismes existe, pour \(\overline{\frR},\overline{\frS}\), l'anneau-produit\footnote{Anneau-produit signifie donc seulement: l'anneau engendré par deux anneaux ayant un sur-anneau commun; le produit n'a pas besoin d'être direct. -- Le fait qu'un anneau-produit n'ait pas toujours à exister pour une intersection donnée est montré par l'exemple suivant: soit \(\frako\) l'anneau des entiers; posons \(\frR=\frako[x]\), \(\frS=\frako[y]\) avec \(2x-1=0\), \(2y-1=0\), de sorte que \(\frako\) soit l'intersection de \(\frR\) et \(\frS\), si aucune combinaison n'est fixée entre \(x\) et \(y\). Mais si \(\frR\) et \(\frS\) se trouvent dans un sur-anneau commun, alors on obtient: \((2x-1)y-x(2y-1)=0\), donc \(x=y\). Il n'existe donc pas d'anneau-produit ayant pour intersection \(\frako\).} \(\overline{\frT}\), lequel, en raison de la permutabilité élément par élément, consiste en tous les éléments de la forme
\[
        \sum_i \bar r_i\bar s_i+\bar r+\bar s
\]
(si \(\overline{\frR},\overline{\frS}\) possèdent des éléments unités, les termes supplémentaires \(\bar r,\bar s\) disparaissent). S'il existe maintenant aussi un anneau-produit \(\frT\), dans lequel \(\frR\) et \(\frS\) sont permutables élément par élément, celui-ci consiste de façon correspondante en tous les éléments de la forme \(\sum_i r_i s_i+r+s\). Les homomorphismes
\[
        \frR\to\overline{\frR},\qquad \frS\to\overline{\frS}
\]
se prolongent donc, par la correspondance
\[
        \sum_i r_i s_i+r+s
        \longmapsto
        \sum_i \bar r_i\bar s_i+\bar r+\bar s
\]
en un homomorphisme \(\frT\to\overline{\frT}\); ainsi la combinaison
\[
        m\Bigl(\sum_i r_i s_i+r+s\Bigr)
        =\sum_i (m r_i)s_i+m r+m s
\]
est univoque et définit, d'après \S 1, 3., un \(\frT\)-module qui englobe le \(\frR,\frS\)-module donné.

\clearpage
